\chapter{The reusable machine}
\label{chap:groups}
\chaptertag{A quick introduction to finite groups}

The necklace proof is vivid, but it was tailored to one symmetry. Group theory isolates the small collection of rules shared by rotations, invertible residues, permutations, and many other systems.

\section{One operation, four rules}

\begin{definition}
A \emph{group} is a set $G$ with an operation $\ast$ such that:
\begin{enumerate}
  \item $a\ast b\in G$ for all $a,b\in G$ (closure);
  \item $(a\ast b)\ast c=a\ast(b\ast c)$ (associativity);
  \item there is an identity $e$ with $e\ast a=a\ast e=a$;
  \item each $a$ has an inverse $a^{-1}$ with $a\ast a^{-1}=a^{-1}\ast a=e$.
\end{enumerate}
If also $a\ast b=b\ast a$, the group is \emph{abelian}.
\end{definition}

\begin{center}
\begin{tabular}{lll}
\toprule
Group & Operation & Identity / inverse \\
\midrule
$\Z/n\Z$ & addition & $0$ / $-a$ \\
$(\Z/n\Z)^\times$ & multiplication of units & $1$ / modular inverse \\
rotations of a regular $n$-gon & composition & no rotation / reverse rotation \\
permutations of a set & composition & do nothing / undo \\
\bottomrule
\end{tabular}
\end{center}

Not every familiar set works: all residues modulo $n$ under multiplication fail because nonunits have no inverse.

\section{Powers and cyclic subgroups}

For $g\in G$, the list
\[
e,g,g^2,g^3,\dots
\]
eventually repeats when $G$ is finite. The first positive $m$ with $g^m=e$ is the \emph{order} of $g$, written $\operatorname{ord}(g)$. The powers form the cyclic subgroup $\angles g$.

\begin{center}
\begin{tikzpicture}
\foreach \k in {0,...,5}{
  \node[circle,draw=teal,fill=mist,minimum size=8mm] (n\k) at ({90-60*\k}:1.65) {$g^{\k}$};
}
\foreach \k [evaluate=\k as \j using {int(mod(\k+1,6))}] in {0,...,5}{
  \draw[-{Stealth},navy!60] (n\k) to[bend left=12] (n\j);
}
\node at (0,0) {multiply by $g$};
\end{tikzpicture}

\small A cyclic group of order $6$: repeated multiplication traces a loop.
\end{center}

\section{Cosets and Lagrange's theorem}

If $H$ is a subgroup of $G$ and $g\in G$, the left coset
\[
gH=\{gh:h\in H\}
\]
is a translated copy of $H$. Two cosets are either equal or disjoint, and every coset has $\abs H$ elements because $h\mapsto gh$ is a bijection.

\begin{theorem}[Lagrange]
If $G$ is finite and $H\le G$, then $\abs H$ divides $\abs G$.
\end{theorem}

\begin{proof}
The distinct cosets of $H$ partition $G$ into equally sized blocks. If there are $r$ blocks, then $\abs G=r\abs H$.
\end{proof}

\begin{corollary}
For every $g$ in a finite group $G$, $\operatorname{ord}(g)\divides\abs G$ and hence $g^{\abs G}=e$.
\end{corollary}

\section{Fermat's theorem, now in one line}

For prime $p$, every nonzero residue has an inverse, so
\[
(\Z/p\Z)^\times=\{1,2,\dots,p-1\}
\]
is a group of order $p-1$. Lagrange immediately yields the standard group-theoretic proof:
\[
a^{p-1}\congr1\pmod p\qquad(p\ndivides a).
\]

The gain is not a shorter proof at any price. The gain is a theorem that can be reused.

\begin{theorem}[Euler]
If $\gcd(a,n)=1$, then
\[
a^{\varphi(n)}\congr1\pmod n,
\]
where $\varphi(n)=\abs{(\Z/n\Z)^\times}$.
\end{theorem}

\begin{proof}
Apply the corollary of Lagrange to the group $(\Z/n\Z)^\times$.
\end{proof}

\begin{bigidea}
The necklace proof uses orbits of a cyclic action; Lagrange uses cosets of a subgroup. Both prove divisibility by partitioning a finite set into equal blocks. Group theory turns that shared move into a portable language.
\end{bigidea}

\section{A second standard proof without naming Lagrange}

If $a\not\congr0\pmod p$, multiplication by $a$ permutes the nonzero residues:
\[
a,2a,\dots,(p-1)a\pmod p
\]
is just $1,2,\dots,p-1$ in another order. Multiply all entries:
\[
a^{p-1}(p-1)!\congr(p-1)!\pmod p.
\]
Since $(p-1)!$ is invertible, cancel it. This is the group proof with the machinery folded away.

\begin{exercisebox}
\begin{enumerate}
  \item List $(\Z/12\Z)^\times$ and make its multiplication table. Find the order of every element. \exlevel{Fluency}
  \item Find $\varphi(20)$ and use Euler's theorem to reduce $3^{2026}$ modulo $20$. \exlevel{Fluency}
  \item Prove that the identity of a group is unique and that each element has a unique inverse. \exlevel{Reasoning}
  \item Prove that two left cosets of $H$ are equal or disjoint. \exlevel{Reasoning}
  \item Show that a group of prime order is cyclic. \exlevel{Reasoning}
  \item Draw the subgroup lattice of $(\Z/12\Z,+)$. Check Lagrange's theorem at every node. \exlevel{Exploration}
  \item Prove $\varphi(mn)=\varphi(m)\varphi(n)$ when $\gcd(m,n)=1$ by using the Chinese remainder theorem. \exlevel{Challenge}
  \item Compare the necklace, permutation, and Lagrange proofs of Fermat's theorem. For each, identify the set being partitioned or permuted and the exact place primality enters. \exlevel{Synthesis}
\end{enumerate}
\end{exercisebox}
